\chapter{CONCLUSION AND FUTURE WORK}
\label{ch:conclusion}

In conclusion, Phase I of the project has produced a complete, tested prediction pipeline for explainable visual commonsense reasoning about social interactions. The system accepts image, video and live-camera input, standardises it with OpenCV, performs visual understanding and language reasoning with SmolVLM, enriches the detected concepts with ConceptNet, and fuses the resulting evidence in a configurable social-interaction reasoning module.

SmolVLM was fine-tuned on the VCR dataset with LoRA for 20 epochs, and the adapter from epoch 11 was selected on a held-out selection set. On separate held-out test questions it improves accuracy over the base model from 44.7\% to 58.7\% for \QtoA, from 41.3\% to 57.7\% for \QAtoR{} and from 15.3\% to 32.7\% for \QtoAR.

The main output of Phase I is a structured prediction of the activity, relationship and intention, accompanied by confidence values, alternatives, evidence and a textual explanation. A reusable \code{Phase1Pipeline} ties all the modules together, the FastAPI--React interface shows the full flow working from start to finish, and all 23 automated tests pass. This meets the Phase I objectives (1 to 10 in Section~\ref{sec:objectives}), and Phase II can build on an implementation that has already been tested.

In Phase II, the focus moves from making predictions to checking the explanations behind them. Grad-CAM maps computed over the SmolVLM vision encoder will give a visual explanation path, and the concepts mentioned in the generated explanation will give a separate textual path. We will measure how well the two concept sets agree using the Jaccard-based cross-modal consistency score (Equation~\ref{eq:cmc}). Faithfulness will be evaluated with the Insertion--Deletion metric and the Pointing Game, and the input frame, the predicted interaction, the three labels, the explanations and the consistency score will be presented in an explainability dashboard.

We also plan the following improvements:
\begin{itemize}
\item building a labelled evaluation set for the activity, relationship and intention predictions of the full pipeline, and running an ablation study on the fusion weights and on each source of evidence;
\item measuring how well the confidence values are calibrated, and improving them;
\item evaluating on a larger VCR test set and training on more VCR questions;
\item batching the VLM prompts to speed up inference, with the aim of getting close to real-time video analysis;
\item adding privacy safeguards so that the system does not name individuals, and auditing its predictions across demographic groups.
\end{itemize}
